\documentclass{article}
\usepackage{amsmath, amssymb, amsthm}
\usepackage{minted}
\usepackage{geometry}
\geometry{left=2cm, right=2cm, top=2cm, bottom=2cm}

\newcommand{\xs}{x_S}
\newcommand{\xt}{x_t}
\newcommand{\taut}{\tau_t}


\title{MATH4907 Assignment 2}
\author{Sanyang Liu}
\date{September 2026}

\begin{document}

\maketitle

\section*{Question 1}

Suppose that \(r(s)\) and \(\sigma(s)>0\) are continuous
functions of time to maturity \(s\). Use the transformation
given in the lecture notes:
\begin{equation*}
    \begin{cases}
        x(S,t)=\ln S+
        \displaystyle\int_0^{T-t}
        \left(r(z)-\frac12\sigma^2(z)\right)\,dz, \\[6pt]
        \tau(t)=\displaystyle\frac12
        \int_0^{T-t}\sigma^2(z)\,dz, \\[6pt]
        V(S,t)=
        \exp\left(-\displaystyle\int_0^{T-t}r(z)\,dz\right)
        u(x,\tau).
    \end{cases}
\end{equation*}
Define
\[
E(t):=\exp\left(-\int_0^{T-t}r(z)\,dz\right).
\]
Then
\[
V(S,t)=E(t)u\bigl(x(S,t),\tau(t)\bigr),
\qquad E'(t)=r(T-t)E(t).
\]
With this time-to-maturity convention, the Black--Scholes
equation in calendar time \(t\) is
\[
V_t+\frac12\sigma^2(T-t)S^2V_{SS}
+r(T-t)SV_S-r(T-t)V=0.
\]
From the definitions of \(x\) and \(\tau\), we compute
\[
\begin{cases}
x_S=\dfrac1S, \\[6pt]
x_{SS}=-\dfrac1{S^2}, \\[6pt]
x_t=\dfrac12\sigma^2(T-t)-r(T-t), \\[6pt]
\tau_t=-\dfrac12\sigma^2(T-t).
\end{cases}
\]
Using the chain rule and \(V=Eu\), we obtain
\[
\begin{cases}
V_S=E\,u_x x_S
    =E\,\dfrac{u_x}{S}, \\[8pt]
V_{SS}
=E\left(u_{xx}x_S^2+u_xx_{SS}\right)
=E\,\dfrac{u_{xx}-u_x}{S^2}, \\[8pt]
V_t=E\left(u_xx_t+u_\tau\tau_t\right)+uE'(t)\\[4pt]
\phantom{V_t}
=E\left[
u_x\left(\dfrac12\sigma^2(T-t)-r(T-t)\right)
-\dfrac12\sigma^2(T-t)u_\tau
+r(T-t)u
\right].
\end{cases}
\]
Substituting these expressions into the Black--Scholes
equation and dividing by \(E>0\), we get
\[
\begin{aligned}
0={}&
u_x\left(\frac12\sigma^2(T-t)-r(T-t)\right)
-\frac12\sigma^2(T-t)u_\tau+r(T-t)u\\
&+\frac12\sigma^2(T-t)(u_{xx}-u_x)
+r(T-t)u_x-r(T-t)u.
\end{aligned}
\]
The \(u_x\) terms and the \(u\) terms cancel. Hence
\[
-\frac12\sigma^2(T-t)u_\tau
+\frac12\sigma^2(T-t)u_{xx}=0.
\]
Since \(\sigma(T-t)>0\), we conclude
\[
\boxed{u_\tau=u_{xx}}
\]
on the domain
\[
(x,\tau)\in\mathbb R\times
\left[0,\frac12\int_0^T\sigma^2(z)\,dz\right].
\]
\qed
\section*{Question 2}
Let
\[
c_{-1}f(x-h)+c_0f(x)+c_1f(x+h)
\]
be the desired finite difference approximation. By Taylor expansion at \(x\),
\begin{align*}
f(x-h)&=f(x)-hf'(x)+\frac{h^2}{2}f''(x)-\frac{h^3}{6}f'''(x)+O(h^4),\\
f(x+h)&=f(x)+hf'(x)+\frac{h^2}{2}f''(x)+\frac{h^3}{6}f'''(x)+O(h^4).
\end{align*}
Thus
\begin{align*}
&c_{-1}f(x-h)+c_0f(x)+c_1f(x+h)\\
&=(c_{-1}+c_0+c_1)f(x)
+h(-c_{-1}+c_1)f'(x)\\
&\quad+\frac{h^2}{2}(c_{-1}+c_1)f''(x)
+\frac{h^3}{6}(-c_{-1}+c_1)f'''(x)+O(h^4).
\end{align*}
To approximate \(h^2f''(x)\), we require the linear system
\[
\begin{cases}
c_{-1}+c_0+c_1=0,\\[2pt]
-c_{-1}+c_1=0,\\[2pt]
\dfrac{1}{2}(c_{-1}+c_1)=1.
\end{cases}
\]
In matrix form,
\[
\begin{pmatrix}
1 & 1 & 1\\
-1 & 0 & 1\\
1 & 0 & 1
\end{pmatrix}
\begin{pmatrix}
c_{-1}\\ c_0\\ c_1
\end{pmatrix}
=
\begin{pmatrix}
0\\ 0\\ 2
\end{pmatrix}.
\]
Solving gives
\[
(c_{-1},c_0,c_1)=(1,-2,1).
\]
Therefore,
\[
f''(x)=\frac{f(x-h)-2f(x)+f(x+h)}{h^2}+O(h^2).
\]
\qed

\section*{Question 3}
We seek coefficients \(c_{++},c_{+-},c_{-+},c_{--}\) such that
\[
\begin{aligned}
S:=\;&c_{++}f(x+h,y+k)+c_{+-}f(x+h,y-k)\\
&+c_{-+}f(x-h,y+k)+c_{--}f(x-h,y-k)
\end{aligned}
\]
approximates \(4hk\,f_{xy}(x,y)\). Assume that \(f\) has
continuous fourth-order partial derivatives near \((x,y)\).\\
Following the hint, let \(a,b\in\{1,-1\}\). The bivariate Taylor
expansion at \((x,y)\) up to third order is
\[
\begin{aligned}
f(x+ah,y+bk)
={}&f+ahf_x+bkf_y\\
&+\frac12\bigl(h^2f_{xx}+2ab\,hkf_{xy}+k^2f_{yy}\bigr)\\
&+\frac16\bigl(
ah^3f_{xxx}+3b\,h^2kf_{xxy}
+3a\,hk^2f_{xyy}+bk^3f_{yyy}
\bigr)+R_{ab}.
\end{aligned}
\]
The four choices of \((a,b)\) give the expansions at the four
points in \(S\). For the constant and first-order terms to
vanish and the coefficient of \(hkf_{xy}\) to equal \(4\), we
require
\[
\begin{cases}
c_{++}+c_{+-}+c_{-+}+c_{--}=0,\\[2pt]
c_{++}+c_{+-}-c_{-+}-c_{--}=0,\\[2pt]
c_{++}-c_{+-}+c_{-+}-c_{--}=0,\\[2pt]
c_{++}-c_{+-}-c_{-+}+c_{--}=4.
\end{cases}
\]
Equivalently,
\[
\begin{pmatrix}
1&1&1&1\\
1&1&-1&-1\\
1&-1&1&-1\\
1&-1&-1&1
\end{pmatrix}
\begin{pmatrix}
c_{++}\\c_{+-}\\c_{-+}\\c_{--}
\end{pmatrix}
=
\begin{pmatrix}
0\\0\\0\\4
\end{pmatrix}.
\]
Solving gives
\[
(c_{++},c_{+-},c_{-+},c_{--})=(1,-1,-1,1).
\]
With these coefficients, the pure second-order terms and all
third-order terms also cancel. Thus
\[
S=f(x+h,y+k)-f(x+h,y-k)
  -f(x-h,y+k)+f(x-h,y-k).
\]
We now estimate the remainder. Define
\[
g(s):=f(x+h,s)-f(x-h,s).
\]
Then \(S=g(y+k)-g(y-k)\). The one-dimensional central Taylor
expansion gives
\[
S=2k\,g'(y)+R_y,\qquad
|R_y|\leq\frac{k^3}{3}
\sup_{|s-y|\leq k}|g'''(s)|.
\]
Since
\[
g'''(s)=f_{yyy}(x+h,s)-f_{yyy}(x-h,s),
\]
the mean value theorem and boundedness of \(f_{xyyy}\) imply
\(g'''(s)=O(h)\). Hence \(R_y=O(hk^3)\).\\
Applying the central Taylor expansion in \(x\) to \(f_y\), with
bounded \(f_{xxxy}\), gives
\[
g'(y)
=f_y(x+h,y)-f_y(x-h,y)
=2h f_{xy}(x,y)+O(h^3).
\]
Consequently,
\[
S=4hk\,f_{xy}(x,y)+O(h^3k+hk^3).
\]
Dividing by \(4hk\), we obtain
\[
\frac{S}{4hk}=f_{xy}(x,y)+O(h^2+k^2).
\]
Since \(O(h^2+k^2)\) implies \(O(h^2+k^2+hk)\), this proves
the stated formula:
\[
\boxed{
f_{xy}(x,y)=
\frac{
f(x+h,y+k)-f(x+h,y-k)
-f(x-h,y+k)+f(x-h,y-k)
}{4hk}
+O(h^2+k^2+hk).
}
\]
\qed
\section*{Question 4}
\subsection*{(a)}

Let
\[
\widetilde{\mathbf u}_j
=(u_{0,j},u_{1,j},\ldots,u_{n,j},u_{n+1,j})^T,
\qquad \lambda=\frac{k}{h^2}.
\]
Following the sign convention for \(B\) in the lecture notes, define
the augmented matrix
\[
\widetilde B=
\begin{pmatrix}
0&0&0&&&\\
-1&2&-1&&&\\
&-1&2&-1&&\\
&&\ddots&\ddots&\ddots&\\
&&&-1&2&-1\\
&&&&0&0
\end{pmatrix}
\in\mathbb R^{(n+2)\times(n+2)}.
\]
Its first and last rows are zero, while each interior row satisfies
\[
(\widetilde B\widetilde{\mathbf u}_j)_i
=-u_{i-1,j}+2u_{i,j}-u_{i+1,j},
\qquad i=1,\ldots,n.
\]
At each interior point, the explicit scheme is
\[
u_{i,j+1}
=u_{i,j}
+\lambda(u_{i-1,j}-2u_{i,j}+u_{i+1,j}),
\qquad i=1,\ldots,n.
\]
The implicit scheme is
\[
u_{i,j+1}
=u_{i,j}
+\lambda(u_{i-1,j+1}-2u_{i,j+1}+u_{i+1,j+1}),
\qquad i=1,\ldots,n.
\]
The endpoint values are
\[
u_{0,j}=\phi_l(\tau_j),\qquad
u_{n+1,j}=\phi_r(\tau_j).
\]
Taking \(\theta=\frac12\) in the theta scheme gives
\[
\begin{aligned}
u_{i,j+1}-u_{i,j}
={}&\frac{\lambda}{2}
 \bigl(u_{i-1,j}-2u_{i,j}+u_{i+1,j}\bigr)\\
&+\frac{\lambda}{2}
 \bigl(u_{i-1,j+1}-2u_{i,j+1}+u_{i+1,j+1}\bigr),
\qquad i=1,\ldots,n.
\end{aligned}
\]
Including the endpoints, its matrix form is
\[
\boxed{
\left(I+\frac{\lambda}{2}\widetilde B\right)
\widetilde{\mathbf u}_{j+1}
=
\left(I-\frac{\lambda}{2}\widetilde B\right)
\widetilde{\mathbf u}_{j}
+\mathbf g_j,
}
\]
where
\[
\mathbf g_j=
\begin{pmatrix}
\phi_l(\tau_{j+1})-\phi_l(\tau_j)\\
0\\
\vdots\\
0\\
\phi_r(\tau_{j+1})-\phi_r(\tau_j)
\end{pmatrix}.
\]
The first and last rows of this matrix equation update the endpoint
values to \(\phi_l(\tau_{j+1})\) and \(\phi_r(\tau_{j+1})\), respectively.

\subsection*{(b)}

Fix \(i\) and \(j\), and write \(x=x_i\) and \(\tau=\tau_j\).
In the definition of \(\eta_{i,j}\), replace
\(v_{i\pm1,j}=u(x\pm h,\tau)\) and
\(v_{i\pm1,j+1}=u(x\pm h,\tau+k)\). Since
\(\lambda=k/h^2\), we have
\[
\begin{aligned}
\eta_{i,j}
={}&u(x,\tau+k)-u(x,\tau)\\
&-\frac{k}{2h^2}
 \Bigl[
 u(x-h,\tau+k)-2u(x,\tau+k)+u(x+h,\tau+k)\\
&\hspace{35mm}
 +u(x-h,\tau)-2u(x,\tau)+u(x+h,\tau)
 \Bigr].
\end{aligned}
\]
All derivatives below are evaluated at \((x,\tau)\), unless their
arguments are displayed. Taylor expansion in time gives
\[
u(x,\tau+k)-u(x,\tau)
=ku_\tau+\frac{k^2}{2}u_{\tau\tau}+O(k^3).
\]
Taylor expansion of the two spatial neighbours about \(x\) gives
\[
u(x-h,\tau)-2u(x,\tau)+u(x+h,\tau)
=h^2u_{xx}+O(h^4).
\]
Applying the same spatial expansion at time \(\tau+k\), then expanding
\(u_{xx}(x,\tau+k)\) about \(\tau\), gives
\[
\begin{aligned}
&u(x-h,\tau+k)-2u(x,\tau+k)+u(x+h,\tau+k)\\
&\qquad=h^2u_{xx}(x,\tau+k)+O(h^4)\\
&\qquad=h^2u_{xx}+kh^2u_{xx\tau}
          +O(k^2h^2+h^4).
\end{aligned}
\]
Thus all the quantities \(u(x\pm h,\tau)\) and \(u(x\pm h,\tau+k)\)
have been expanded about the fixed point \((x_i,\tau_j)\).
Substituting the expansions into \(\eta_{i,j}\) yields
\[
\begin{aligned}
\eta_{i,j}
&=ku_\tau+\frac{k^2}{2}u_{\tau\tau}
 -\frac{k}{2h^2}
  \left(2h^2u_{xx}+kh^2u_{xx\tau}\right)
 +O(k^3+kh^2)\\
&=k(u_\tau-u_{xx})
 +\frac{k^2}{2}(u_{\tau\tau}-u_{xx\tau})
 +O(k^3+kh^2).
\end{aligned}
\]
The heat equation implies \(u_\tau=u_{xx}\), and differentiating
it with respect to \(\tau\) gives
\(u_{\tau\tau}=u_{xx\tau}\). Hence the displayed leading terms
cancel, and
\[
|\eta_{i,j}|\le C(kh^2+k^3).
\]
By Young's inequality with conjugate exponents
\(p=3/2\) and \(q=3\) (i.e., $\frac{1}{p}+\frac{1}{q}=1$),
\[
kh^2
\le \frac{(h^2)^{3/2}}{3/2}+\frac{k^3}{3}
=\frac23h^3+\frac13k^3.
\]
Consequently,
\[
|\eta_{i,j}|
\le C\left(\frac23h^3+\frac43k^3\right)
\le \frac{4C}{3}(h^3+k^3).
\]
Hence we take
\[
\boxed{\alpha=3,\qquad\beta=3.}
\]

\section*{Question 5}

Let \(x=\ln S\) and \(\tau=T-t\). Write \(V(x,\tau)\) for the
European call price. With
\[
S_0=K=100,\qquad r=0.05,\qquad \sigma=0.15,\qquad T=1,
\]
the pricing equation is
\[
V_\tau
=\frac{\sigma^2}{2}V_{xx}
+\left(r-\frac{\sigma^2}{2}\right)V_x-rV.
\]
Define
\[
\nu:=r-\frac{\sigma^2}{2}=0.03875.
\]
The initial condition at \(\tau=0\) is the call payoff
\[
V(x,0)=(e^x-K)^+.
\]

\subsection*{(a)}

Let \(x_i=x_{\min}+ih\), \(\tau_j=jk\), and
\(V_{i,j}\approx V(x_i,\tau_j)\). Applying a forward difference
in \(\tau\) and second-order central differences in \(x\) gives
\[
\frac{V_{i,j+1}-V_{i,j}}{k}
=
\frac{\sigma^2}{2}
\frac{V_{i-1,j}-2V_{i,j}+V_{i+1,j}}{h^2}
+\nu\frac{V_{i+1,j}-V_{i-1,j}}{2h}
-rV_{i,j}.
\]
Hence, at each interior point,
\[
\boxed{
V_{i,j+1}
=c_-V_{i-1,j}+c_0V_{i,j}+c_+V_{i+1,j},
}
\]
where
\[
c_-=\frac{\sigma^2k}{2h^2}-\frac{\nu k}{2h},
\qquad
c_0=1-\frac{\sigma^2k}{h^2}-rk,
\qquad
c_+=\frac{\sigma^2k}{2h^2}+\frac{\nu k}{2h}.
\]

\subsection*{(b)}

On the truncated interval \([-1,5]\), we use the usual
low- and high-stock-price boundary approximations for a
European call:
\[
\boxed{
V(-1,\tau)=0,\qquad
V(5,\tau)=e^5-Ke^{-r\tau}
        =e^5-100e^{-0.05\tau},
\quad 0\leq\tau\leq1.
}
\]
The first boundary approximates the call value by zero when
\(S=e^{-1}\) is small. The second uses the large-\(S\) behaviour
\(V\approx S-Ke^{-r\tau}\).

\subsection*{(c)}
Use \(h=0.01\) and \(k=1/225\). Starting from
\[
V_{i,0}=(e^{x_i}-100)^+,
\]
apply the explicit update in part (a) for \(225\) time steps.
The following Python code also uses linear interpolation to
evaluate the result at \(x=\ln100\), which is not a grid point. See the code in the Appendix. For the interval \([-1,5]\), the computed result is
\[
\boxed{V(\ln100,1)\approx 8.59888676.}
\]

\subsection*{(d)}

Changing the spatial interval to \([-2,10]\) while keeping
\(h=0.01\) and \(k=1/225\), the boundary approximations become
\[
V(-2,\tau)=0,\qquad
V(10,\tau)=e^{10}-100e^{-0.05\tau}.
\]
Running the same code with
\texttt{option\_value(-2.0, 10.0)} gives
\[
\boxed{V(\ln100,1)\approx 8.59888870.}
\]
The difference between the results in parts (c) and (d) is
approximately \(1.94\times10^{-6}\).

\newpage
\section*{Appendix: Code for Q5(c)}
\small
\begin{minted}{python}
import numpy as np

# Model parameters
K = 100.0          # Strike price
S0 = 100.0         # Stock price at which we evaluate the option
r = 0.05           # Risk-free interest rate
sigma = 0.15       # Volatility

# Discretization
h = 0.01           # Step size in x = ln(S)
m = 225            # Number of time steps
k = 1.0 / m        # Time step size
nu = r - sigma**2 / 2

# Coefficients of the explicit finite difference scheme:
# V_i^{j+1} = c_minus V_{i-1}^j + c_zero V_i^j + c_plus V_{i+1}^j
c_minus = sigma**2 * k / (2 * h**2) - nu * k / (2 * h)
c_zero = 1 - sigma**2 * k / h**2 - r * k
c_plus = sigma**2 * k / (2 * h**2) + nu * k / (2 * h)

def value(a, b):
    # Construct the spatial grid on [a, b]
    n = round((b - a) / h)
    x = a + h * np.arange(n + 1)

    # At tau = 0, the option value equals its payoff
    V = np.maximum(np.exp(x) - K, 0.0)
    new_V = np.empty_like(V)

    # Advance the solution from tau = 0 to tau = 1
    for j in range(m):
        # Update all interior grid points using values at the old time level
        new_V[1:-1] = (
            c_minus * V[:-2]
            + c_zero * V[1:-1]
            + c_plus * V[2:]
        )

        # Set boundary values at the new time level
        tau_next = (j + 1) * k
        new_V[0] = 0.0
        new_V[-1] = np.exp(b) - K * np.exp(-r * tau_next)

        # Swap the old and new arrays for the next time step
        V, new_V = new_V, V

    # ln(S0) is generally between grid points: interpolate linearly
    target = np.log(S0)
    position = (target - a) / h
    i = int(np.floor(position))
    weight = position - i

    return (1.0 - weight) * V[i] + weight * V[i + 1]

print("[-1, 5]: ", value(-1.0, 5.0))
print("[-2, 10]:", value(-2.0, 10.0))
\end{minted}

\end{document}
